\section{Conclusion and Future Work}
\label{sec:conclusion}
This article has presented \method, an efficient 3D vision-language-action model built upon a pre-trained vision-language model (VLM)~\cite{beyer2024paligemma}, together with its memory-augmented extension \memmethod.
\method rests on a single alignment principle: 3D observations are rendered as multi-view 2D images to match the input space of the VLM, actions are expressed as 2D heatmaps in the same image space, and a scalable pre-training stage teaches the VLM to predict heatmaps before it is fine-tuned for action prediction.
\memmethod extends this principle with a unified spatio-temporal memory: temporal memory retains the interaction history to determine \emph{what to do next}, while spatial memory re-renders previously observed geometry to determine \emph{where exactly to act}.
Extensive experiments on standard and memory-dependent benchmarks, in both simulation and the real world, show that the framework learns 3D manipulation efficiently and effectively, and that the memory is strictly additive to the base policy.
Future work includes broadening pre-training to more diverse tasks such as semantic segmentation and keypoint detection, adopting more expressive action decoders, and replacing the annotation-dependent sub-goal gate with self-supervised alternatives.
Moreover, since the token-space memory injection is agnostic to timescale, 
extending this mechanism to support cross-episode or life-long memory represents a promising path toward robots that continuously accumulate, refine, and transfer manipulation skills across long operational horizons.